%%%%%%%%%%%%%%%%%%%%%%%%%%%%%%%%%%%%%%%
% Cover Letter - FIX, Менеджер проекта (IT, B2B/B2G)
%%%%%%%%%%%%%%%%%%%%%%%%%%%%%%%%%%%%%%%

\documentclass[]{cover}
\usepackage{fancyhdr}

\pagestyle{fancy}
\fancyhf{}

\rfoot{Page \thepage \hspace{0pt}}
\thispagestyle{empty}
\renewcommand{\headrulewidth}{0pt}

\begin{document}

%%%%%%%%%%%%%%%%%%%%%%%%%%%%%%%%%%%%%%
%     TITLE NAME
%%%%%%%%%%%%%%%%%%%%%%%%%%%%%%%%%%%%%%
\namesection{}{\Huge{Алексей Орлов}}{  \href{mailto:lex.orlov78@yandex.ru}{lex.orlov78@yandex.ru} | +7~(901)~797-54-88 |  \urlstyle{same}\href{https://www.linkedin.com/in/alex-orlov-v/}{LinkedIn}
}

%%%%%%%%%%%%%%%%%%%%%%%%%%%%%%%%%%%%%%
%     MAIN COVER LETTER CONTENT
%%%%%%%%%%%%%%%%%%%%%%%%%%%%%%%%%%%%%%

\currentdate{16 сентября 2026 г.}
\lettercontent{Здравствуйте, команда FIX,}

\lettercontent{Меня заинтересовала позиция менеджера проекта на продукте по противодействию нелегальной активности в интернете. Последние 6 лет я решал именно ту задачу, которая стоит в центре этой роли: переводить приоритеты заказчика в измеримую дорожную карту и держать бюджет, сроки и команду под контролем, - сначала в Банке ВТБ, затем в ООО «ПочтаТех».}

\lettercontent{FIX - продуктовая, а не аутсорс-компания, и это принципиально: рост от казанского стартапа до международного холдинга с офисами в Москве, Казани, Лондоне, Дубае и Лимасоле говорит о зрелой продуктовой культуре, а не о разовых проектах на заказ. Отдельно откликается то, что компания компенсирует использование ИИ-агентов сотрудникам - я не просто умею работать с ИИ-инструментами, а системно встраиваю их в управление проектом.}

\lettercontent{Что я принесу на позицию менеджера проекта:}

{\raggedright\fontspec{DejaVu Sans}\fontsize{11pt}{13pt}\selectfont
\begin{itemize}
    \item \textbf{Дорожная карта и приоритизация бэклога:} спроектировал модель квартального планирования и приоритизации инициатив на основе аналитики и бизнес-целей в ООО «ПочтаТех», сократившую Time-to-Market примерно на 20\%
    \item \textbf{Бюджет и отчётность:} вёл план-факт по срокам и стоимости портфеля проектов, настроил систему метрик (TTM, Cycle Time, SLA) для регулярной отчётности перед CTO и руководителями направлений
    \item \textbf{Управление технической командой:} руководил группой разработчиков (Все Инструменты.РУ, Oracle BI/ETL) и на протяжении карьеры координировал работу с разработчиками и аналитиками
    \item \textbf{Переговоры и согласования:} сертификация ICF-коуча и опыт фасилитации PI Planning в Банке ВТБ (до 15 команд, 70+ участников) - инструменты для сложных многосторонних обсуждений
\end{itemize}\par}
\vspace{6pt}

\lettercontent{Честно: мой опыт управления заказчиком - это внутренние стейкхолдеры (CTO, руководители направлений, бизнес-подразделения), а не внешние коммерческие B2B/B2G-клиенты, которых требует вакансия. Но задача по сути та же - быть ключевым звеном между заказчиком и командой, - и техническое образование (к.т.н.) вместе с опытом ведения бюджета и дорожной карты дают быстрый старт в этой роли. Буду рад обсудить, как мой опыт может усилить команду продукта.}

\begin{flushright}
\closing{С уважением,}

\signature{Алексей Орлов}
\end{flushright}
\end{document}
